\unitopen{02}
\begin{candos}
\item powiedzieć, jak się nazywasz, skąd jesteś, gdzie mieszkasz i czym się zajmujesz;
\item odmienić regularny czasownik w czasie teraźniejszym (\de{Präsens}) i poprawnie użyć \de{sein};
\item zadać pytanie otwarte (\de{wer, was, wie, wo, woher}) i pytanie zamknięte (tak/nie);
\item poprosić o powtórzenie i wolniejsze mówienie.
\end{candos}

\section{Słownictwo}
\vocabtable{02}

\section{Zaimki osobowe i czasownik \de{sein}}
\begin{center}\small
\begin{tabular}{@{}lll@{}}
\toprule
\textsf{\textbf{Zaimek}} & \textsf{\textbf{\de{sein} (być)}} & \textsf{\textbf{\de{kommen} (przychodzić)}}\\
\midrule
\de{ich} & \de{bin} & \de{komme}\\
\de{du} & \de{bist} & \de{kommst}\\
\de{er / sie / es} & \de{ist} & \de{kommt}\\
\de{wir} & \de{sind} & \de{kommen}\\
\de{ihr} & \de{seid} & \de{kommt}\\
\de{sie / Sie} & \de{sind} & \de{kommen}\\
\bottomrule
\end{tabular}
\end{center}

\de{sein} jest nieregularny i musisz znać go na pamięć. Czasowniki regularne mają końcówki
\textbf{-e, -st, -t, -en, -t, -en} dodawane do rdzenia: \de{komm-}, \de{wohn-}, \de{lern-}.

\subsection*{Wyjątki od końcówek}
Rdzeń na \de{-t} lub \de{-d} wymaga \de{-e-} dla wygody wymowy: \de{arbeiten} $\rightarrow$ \de{du arbeitest},
\de{er arbeitet}. Rdzeń na \de{-ß} lub \de{-s}: \de{heißen} $\rightarrow$ \de{du heißt} (jedno \de{s}).

\begin{zasada}[Czasownik zawsze na drugim miejscu]
W zdaniu oznajmującym odmieniony czasownik stoi na \emph{drugiej pozycji}: \de{Ich komme aus Polen.}
\de{Aus Polen komme ich.} (podmiot po czasowniku, bo na pierwszym miejscu jest inna część zdania).
W polskim szyk jest swobodny; w niemieckim ta reguła jest sztywna.
\end{zasada}

\section{Pytania}
Pytanie otwarte zaczyna się od słowa pytającego, czasownik jest na drugim miejscu:
\de{Woher kommst du?} \de{Wo wohnen Sie?} \de{Wie heißt du?} \de{Was bist du von Beruf?}
Pytanie zamknięte (tak/nie) zaczyna się od \emph{czasownika}: \de{Sprechen Sie Deutsch?} \de{Bist du Pole?}

\begin{falle}
Po polsku mówisz „Jestem inżynierem” (narzędnik). Po niemiecku zawód nie ma rodzajnika i nie ma
przypadku narzędnego: \de{Ich bin Ingenieur.} Nie: \de{Ich bin ein Ingenieur} (tak mówi się dopiero,
gdy dodajesz przymiotnik: \de{Ich bin ein guter Ingenieur}). Druga pułapka: \de{Ich bin 30} jest błędem;
wiek zaczyna się w jednostce~3.
\end{falle}

\begin{szwajcaria}
Po niemiecku piszesz \de{die Schweiz} (z rodzajnikiem), po polsku „Szwajcaria”. Szwajcarzy
rozmawiają o pracy bez zażenowania; \de{Was machen Sie beruflich?} to normalne pytanie w pierwszej minucie
rozmowy. Osoba z Zurychu to \de{Zürcher/in}, ale o narodowości mówisz \de{Schweizer/in}.
\end{szwajcaria}

\begin{praca}
Na standupie: \de{Ich bin Entwickler im Backend-Team.} \de{Ich arbeite an der API.} Sam rzeczownik
\de{Entwickler} pisz wielką literą, jak każdy rzeczownik.
\end{praca}

\section{Dialogi}
\textbf{Dialog 1 --- konferencja (formalnie)}
\begin{dialog}
\turn{A}{Guten Tag. Mein Name ist Anna Keller. Wie heißen Sie?}{Dzień dobry. Nazywam się Anna Keller. Jak Pan się nazywa?}
\turn{B}{Ich heiße Piotr Nowak. Freut mich.}{Nazywam się Piotr Nowak. Miło mi.}
\turn{A}{Freut mich auch. Woher kommen Sie?}{Mnie też. Skąd Pan pochodzi?}
\turn{B}{Ich komme aus Polen. Und Sie?}{Pochodzę z Polski. A Pani?}
\turn{A}{Ich komme aus Zürich. Was sind Sie von Beruf?}{Jestem z Zurychu. Czym się Pan zajmuje?}
\turn{B}{Ich bin Softwareentwickler. Sprechen Sie Englisch?}{Jestem programistą. Mówi Pani po angielsku?}
\turn{A}{Ja, aber wir sprechen heute Deutsch.}{Tak, ale dziś mówimy po niemiecku.}
\turn{B}{Ein bisschen. Langsamer, bitte.}{Trochę. Wolniej, proszę.}
\end{dialog}

\section{Ćwiczenia}
\begin{cwiczenia}
\exitem{Uzupełnij formą \de{sein}: \de{Ich \dots{} Entwickler.}}{\de{bin}}
\exitem{Uzupełnij: \de{Du \dots{} aus Polen.}}{\de{bist}}
\exitem{Uzupełnij: \de{Wir \dots{} Kollegen.}}{\de{sind}}
\exitem{Uzupełnij: \de{Ihr \dots{} in Zürich.}}{\de{seid}}
\exitem{Odmień \de{wohnen} dla \de{er}.}{\de{er wohnt}}
\exitem{Odmień \de{arbeiten} dla \de{du}.}{\de{du arbeitest}}
\exitem{Odmień \de{heißen} dla \de{du}.}{\de{du heißt}}
\exitem{Zadaj pytanie o kraj pochodzenia do kogoś, kogo tykasz.}{\de{Woher kommst du?}}
\exitem{Zadaj pytanie tak/nie o znajomość niemieckiego (formalnie).}{\de{Sprechen Sie Deutsch?}}
\exitem{Popraw błąd: \de{Ich komme Polen.}}{\de{Ich komme aus Polen.}}
\exitem{Ułóż zdanie, zaczynając od \de{Aus Polen}: (ich / komme)}{\de{Aus Polen komme ich.}}
\exitem{Przetłumacz: „Jestem programistą i mieszkam w Zurychu.”}{\de{Ich bin Entwickler und ich wohne in Zürich.}}
\end{cwiczenia}

\section{Trening słownictwa (generowany)}
\drill{02}

\weekend{Napisz i powiedz na głos pięć zdań o sobie (imię, kraj, miasto, zawód, języki), potem
zadaj te same pytania wymyślonej osobie. Nagraj się i porównaj z modelem: \emph{Ich heiße \dots{} Ich komme aus \dots{}
Ich wohne in \dots{} Ich bin \dots{} Ich spreche \dots}.}
